\documentclass{article}

% --- FUENTES Y CODIFICACIÓN ---
\usepackage[final,babel=True]{microtype}
\microtypecontext{spacing=nonfrench}
\usepackage{fontspec}
\usepackage{svg}
\usepackage{afterpage}
\svgsetup{inkscapelatex=false,inkscapearea=drawing}
\usepackage{unicode-math}
\setmathfont{Libertinus Math}[Scale=MatchLowercase]
\setsansfont[Scale=MatchLowercase,Ligatures=TeX,
  Extension=.otf,UprightFont=*-Regular,ItalicFont=*-Italic,
  BoldFont=*-SemiBold,BoldItalicFont=*-SemiBoldItalic]{IBMPlexSansCondensed}
\setmonofont[Scale=0.91,Extension=.otf,UprightFont=*-Regular,
  ItalicFont=IBMPlexMono-Italic.otf,BoldFont=IBMPlexMono-Bold.otf,
  BoldItalicFont=IBMPlexMono-BoldItalic.otf]{IBMPlexMono.otf}

% --- IDIOMAS ---
\usepackage[english,portuguese,main=spanish]{babel}
\usepackage[autostyle=true]{csquotes}
\setmainfont[Numbers={OldStyle,Proportional},
  Ligatures={TeX,Common,Discretionary},Scale=1.18]{Libertinus Serif}
\renewcommand{\normalsize}{\fontsize{10pt}{13pt}\selectfont}

% --- BIBLIOGRAFÍA ---
\input{/home/alberto/.gbpublisher/latex/cita-vancouver-config.tex}


% --- PÁGINA ---
\usepackage[paperwidth=210mm,paperheight=297mm,
  top=13mm,bottom=23mm,left=13mm,right=87mm,
  headsep=0mm,footskip=12mm,headheight=29mm,headsep=10mm,
  marginparsep=6mm,marginparwidth=68mm,
  includehead,includefoot]{geometry}
\setcounter{secnumdepth}{0}

% --- POSICIONAMIENTO ABSOLUTO ---
\usepackage[absolute,overlay]{textpos}
\setlength{\TPHorizModule}{1mm}
\setlength{\TPVertModule}{1mm}
\textblockorigin{0mm}{0mm}
\newlength{\gbbloquex}
\AtBeginDocument{\setlength{\gbbloquex}{\dimexpr 1in+\hoffset+\oddsidemargin+\textwidth+\marginparsep\relax}}

% --- TABLAS Y FIGURAS ---
\usepackage{booktabs}
\usepackage{graphicx}
\usepackage[labelfont=bf,font=small,labelsep=period,format=plain]{caption}

% --- VARIOS ---
\newcommand{\epigraph}[2]{%
  \par\nobreak\noindent\par\nobreak\vspace{.5\baselineskip}%
  \hfill{%
    \small
    \begin{tabular}{@{}>{{\raggedright\arraybackslash}}m{.65\textwidth}@{}}%
      #1 \\[1ex]%
      \midrule%
      \hfill #2%
    \end{tabular}%
  }%
  \vspace{.5\baselineskip}%
}
\usepackage{hyperref}
\hypersetup{colorlinks=True,linkcolor=azulrevista,citecolor=azulrevista,urlcolor=azulrevista,filecolor=azulrevista,
linktoc=all,breaklinks=True,bookmarksopen=True,bookmarksnumbered=True}
\usepackage[bottom,stable,hang]{footmisc}
\usepackage{scrextend}
\usepackage{changepage}
\newcommand{\fullwidthblock}[1]{%
  \begin{adjustwidth}{0pt}{-\dimexpr\marginparsep+\marginparwidth\relax}%
  \centering #1%
  \end{adjustwidth}%
}
\usepackage{rotating}
\setlength{\footnotesep}{10pt}
\renewcommand{\footnoterule}{%
  \kern -3pt%
  \hrule height 0.5pt width 0.4\columnwidth%
  \kern 6pt}
\renewcommand*{\thefootnote}{\scriptsize\sf{[\arabic{footnote}]}}
\makeatletter
\patchcmd{\@footnotetext}{\footnotesize}{\small}{}{}
\makeatother
\newcommand*\footnotemarkspace{0em}
\deffootnote{\footnotemarkspace}
  {\parindent}
  {\makebox[\footnotemarkspace][r]{\llap{\thefootnotemark\quad}}}
\raggedbottom
\clubpenalty=10000
\widowpenalty=10000
\usepackage{siunitx}
\usepackage{enumitem}
\setlist{nosep,topsep=4pt}
\usepackage{titlesec}
\titleformat{\section}[hang]
  {\sf\bfseries\Large\color{azulrevista}\raggedright}
  {\thesection}{.5em}{}[]
\titlespacing{\section}
  {\parindent}{18pt plus 0.5pt minus 0.5pt}{6.75pt}
\titleformat{\subsection}[hang]
  {\sf\large\color{azulrevista}\raggedright}
  {\thesubsection}{.5em}{}[]
\titlespacing{\subsection}
  {\parindent}{18pt plus 0.5pt minus 0.5pt}{6.75pt}
\titleformat{\subsubsection}[hang]
  {\rm\bfseries\normalsize\color{azulrevista}\raggedright}
  {\thesubsubsection}{.5em}{}[]
\titlespacing{\subsubsection}
  {\parindent}{18pt plus 0.5pt minus 0.5pt}{6.75pt}
\titleformat{\paragraph}[runin]
  {\bfseries\color{azulrevista}}{}{0em}{}[\mbox{ --- }]
\titlespacing{\paragraph}
  {0pt}{6.75pt plus 0.5pt minus 0.5pt}{0pt}
\titleformat{\subparagraph}[hang]
  {\sf\bfseries\large\color{azulrevista}\centering}
  {\thesubparagraph}{.5em}{}[]
\titlespacing{\subparagraph}
  {\parindent}{18pt plus 0.5pt minus 0.5pt}{6.75pt}
% --- COLOR ---
\usepackage[svgnames]{xcolor}
\usepackage{tcolorbox}
\definecolor{azulrevista}{HTML}{2d5a8e}

\usepackage{listings}
\lstset{
  basicstyle=\ttfamily\small,
  breaklines=true,
  breakatwhitespace=false,
  frame=single,
  framerule=0pt,
  backgroundcolor=\color{black!5!white},
  fillcolor=\color{black!5!white},
  rulecolor=\color{black!5!white},
  xleftmargin=8pt,xrightmargin=8pt,
  framexleftmargin=8pt,framexrightmargin=8pt,
  framextopmargin=6pt,framexbottommargin=6pt,
  aboveskip=8pt,belowskip=8pt,
  keywordstyle=\color{blue!70!black}\bfseries,
  commentstyle=\color{green!50!black}\itshape,
  stringstyle=\color{orange!80!black},
  showstringspaces=false,
  numbers=none,
  captionpos=b
}
\renewenvironment{quote}{\normalsize\list{}{\sf\leftmargin=14pt\rightmargin=0pt}\item\relax}{\endlist}


\renewenvironment{verse}
  {\normalsize\list{}{\leftmargin=14pt\rightmargin=0pt\parsep=0pt}
   \item\relax\itshape}
  {\endlist}
\newcommand{\articulotipo}{}

% --- HEADER / FOOTER ---
\usepackage{fancyhdr}
\setlength{\headwidth}{\textwidth}
\addtolength{\headwidth}{\marginparsep}
\addtolength{\headwidth}{\marginparwidth}

\fancypagestyle{firstpage}{%
  \fancyhf{}%
  \fancyhead[L]{%
    \begin{minipage}[t]{\headwidth}%
    \includegraphics[width=184mm,height=20mm,keepaspectratio=false]{/home/alberto/00_Produccion/99_Demos/Cicatriz-AR/media/logo-revista.png}\\[3mm]%
      {\setlength{\fboxsep}{0pt}%
      \colorbox{azulrevista}{%
        \parbox[c][6mm][c]{\headwidth}{%
          \hspace{12pt}\sffamily\bfseries\color{white}\articulotipo%
        }%
      }}%
    \end{minipage}%
  }%
  \renewcommand{\headrulewidth}{0pt}%
  \fancyfoot[L]{%
    \begin{minipage}[t]{\headwidth}%
      \textcolor{azulrevista}{\rule{\headwidth}{1pt}}\\[1pt]%
      \sffamily\thepage\hfill%
      \small\articulorevista{} / e-ISSN: \texttt{\articuloissne}{} / Vol.~\articulovolumen{},~n.\textsuperscript{o}~\articulonumero{} / \articulofecha%
    \end{minipage}%
  }%
  \renewcommand{\footrulewidth}{0pt}%
}%

\pagestyle{fancy}
\fancyhf{}
\fancyhead[L]{%
  \raisebox{\dimexpr\headheight-7mm\relax}[0pt][0pt]{%
    \begin{minipage}[b]{\headwidth}%
      \sffamily\thepage\hfill%
      \small\articulorevista{} / e-ISSN: \texttt{\articuloissne}{} / Vol.~\articulovolumen{},~n.\textsuperscript{o}~\articulonumero{} / \articulofecha\\[1pt]%
      \textcolor{azulrevista}{\rule{\headwidth}{1pt}}%
    \end{minipage}%
  }%
}%
\renewcommand{\headrulewidth}{0pt}

% --- MACROS DEL ARTÍCULO (SOBREESCRITAS POR EL XSLT) ---
\newcommand{\articulotitulo}{}
\newcommand{\articulosubtitulo}{}
\newcommand{\articulotitulotrans}{}
\newcommand{\articuloautores}{}
\newcommand{\articulodoi}{}
\newcommand{\articuloidiomaxml}{es}
\newcommand{\articulorevista}{}
\newcommand{\articuloissne}{}
\newcommand{\articulovolumen}{}
\newcommand{\articulonumero}{}
\newcommand{\articulofecha}{}
\newcommand{\articulopagina}{0}

\begin{document}

\selectlanguage{spanish}

\renewcommand{\articuloidiomaxml}{es}
\renewcommand{\articulotitulo}{Editorial}
\renewcommand{\articuloautores}{Érica Lorena Witman}
\renewcommand{\articulotipo}{EDITORIAL}
\renewcommand{\articulorevista}{Cicatriz-AR}
\renewcommand{\articuloissne}{2618-3536}
\renewcommand{\articulovolumen}{14}
\renewcommand{\articulonumero}{1}
\renewcommand{\articulofecha}{ 2026}
\renewcommand{\articulopagina}{1}

\setcounter{page}{1}
\thispagestyle{firstpage}

% --- TÍTULO A ANCHO COMPLETO ---
\noindent\begin{minipage}[t]{\dimexpr\textwidth+\marginparsep+\marginparwidth\relax}%
{\sffamily\bfseries\Large\articulotitulo\par}%
\end{minipage}\par
\vspace{3mm}%
\renewcommand{\articulotitulotrans}{Editor's note}
\noindent\begin{minipage}[t]{\dimexpr\textwidth+\marginparsep+\marginparwidth\relax}%
{\sffamily\itshape\normalsize\articulotitulotrans\par}%
\end{minipage}\par
\vspace{4mm}%

% --- BLOQUE LATERAL DE METADATOS ---
\begin{textblock*}{68mm}(\gbbloquex,92mm)
{\setlength{\leftskip}{0pt}\setlength{\parindent}{0pt}%
\noindent\begin{minipage}[t]{68mm}
\sffamily\small\setlength{\parskip}{1pt}
{\bfseries\color{azulrevista}Érica Lorena Witman}\par
{\scriptsize\texttt{\href{https://orcid.org/0000-0002-7765-2088}{https://orcid.org/0000-0002-7765-2088}}}\par
\vspace{6pt}
{\bfseries\color{azulrevista!70}Derechos de autor\'ia:}\par
\textcopyright\ 2026 Witman, É. Publicado por Cicatriz-AR. Se permite el uso, distribución y reproducción sin fines comerciales, siempre que se cite la obra original y las obras derivadas se distribuyan bajo la misma licencia. (\href{https://creativecommons.org/licenses/by-nc-sa/4.0/}{https://creativecommons.org/licenses/by-nc-sa/4.0/})\par\vspace{6pt}
{\bfseries\color{azulrevista!70}Publicación:} 2026, vol.~14, n\'um.~1, pp.~1--1\par\vspace{6pt}
{\bfseries\color{azulrevista!70}Licencia:} CC BY-NC-SA 4.0\par
{\scriptsize\texttt{\href{https://creativecommons.org/licenses/by-nc-sa/4.0/}{https://creativecommons.org/licenses/by-nc-sa/4.0/}}}\par
\vspace{6pt}
\end{minipage}%
}% FIN GRUPO leftskip
\end{textblock*}


\bigskip
\noindent\rule{\linewidth}{0.4pt}
\bigskip

Para esta nueva edición elegimos «Ciencia y Caridad», el óleo sobre lienzo que Pablo Picasso a principios de 1897 regaló al mundo, a escala real, y se encuentra actualmente en el Museu Picasso en Barcelona, España. Se enmarca en el realismo social, en el subgénero «La pintura hospitalaria»; la paleta artística que el pintor brindó fueron 16 colores para tan solo 16 años de edad de Picasso, obteniendo una mención honorífica en la Exposición Universal de Bellas Artes de Madrid.

El médico, dibujado sentado mientras toma el pulso a la mujer moribunda observando su reloj de bolsillo, encarna la Ciencia y la monja, cuya presencia en los hospitales solía ser habitual, representa a la Caridad. La unión simbólica de las mismas –que origina el título– resulta un choque de etapas de la historia de la medicina y de la enfermería, y motiva a reflexionar sobre la necesidad de unir tanto lo científico como lo espiritual y afectivo. Hoy en día, el término de caridad resulta extraño, y se prefiere el de «justicia social» (igualdad de oportunidades).

No hay que olvidar que la Etapa Vocacional, que es la segunda etapa de la historia de la Enfermería, marcó el ejercicio de las prácticas cuidadoras. En muchos países, el cuidado de los enfermos estaba asociado a las monjas, con más formación y un voto religioso que les impulsaba a cuidar de los más débiles.

Esta preocupación social dio como resultado una serie de cambios, los cuales marcaron el inicio de la enfermería moderna y, con ello, permitieron avanzar en la profesionalización de la actividad del cuidado.

La atención centrada en la persona, encuadre actual, se deja ver a trasluz de las pinceladas, con la mujer pálida y sin fuerzas en el centro del óleo.

A pesar de las dificultades que atravesamos como equipo de salud, la generosidad, vocación de ayuda, ética, humanismo, entusiasmo, creatividad, cuidados, en pos del cierre de las heridas debe brotar de nuestras manos y como bien dice el Dr. Gregorio Marañón, médico endocrinólogo español, «tienen la conciencia cierta de que hasta donde no llega el saber, llega siempre el amor».

Bienvenidos a una sabia y amorosa edición de la \emph{Revista Cicatriz-AR}.


\end{document}
